\section{Practical and theoretical rationality: the scalar-order implication}
\label{sec:p11-rat22-numerical-order-annex}

This item supplies the elementary numerical implication used in Wedgwood's
survey \parencite[p.~140]{RalphWedgwood2021}. It does not supply a
representation-existence theorem, Savage's full theory, or a justification
of the normative premise that rational preferences admit the representation.

\begin{annexitem}{annex:p11.rat22.numerical-order}{thm:p11-rat22-numerical-order}{Exact scalar comparison implies the stated ordering conditions}
\begin{proof}
For any relevant prospects \(A,B\), the real numbers \(U(A),U(B)\)
satisfy exactly one of \(U(A)>U(B)\), \(U(B)>U(A)\), or
\(U(A)=U(B)\). The two assumed biconditionals translate these possibilities
respectively into strict preference for \(A\), strict preference for \(B\),
and indifference. This proves exhaustiveness and exclusivity. When
\(A=B\), equality gives indifference and excludes both strict comparisons.

If \(A\) is strictly preferred to \(B\), then \(U(A)>U(B)\). If \(B\)
is strictly preferred to \(C\), then \(U(B)>U(C)\), and real-order
transitivity gives \(U(A)>U(C)\). If instead \(B\) and \(C\) are
indifferent, then \(U(B)=U(C)\), and substitution gives the same strict
inequality \(U(A)>U(C)\). The representation biconditional therefore
gives strict preference for \(A\) over \(C\) in both cases.

The proof concerns arbitrary individual prospects and never enumerates
the prospect domain. It therefore applies to finite, infinite, singleton,
and empty domains; the universally quantified conclusions are vacuous
in the empty case. No probability, expected-utility, or psychological
assumption is used.
\end{proof}
\end{annexitem}
